\paragraph{Where $\bar H$ came from.} We arrived at this operator through a
rate-distortion analysis of merging, which treats the merge as lossy
compression of $T$ task vectors into one weight vector and asks what worst-case
distortion is achievable at a given storage budget. Appendix~\ref{app:rd}
states the resulting lower bound and a matching construction, and
Appendix~\ref{app:proofs} proves them. We report the analysis as provenance
rather than as a contribution, for three reasons given in full there and
summarised here: its floor term is exactly zero on all four base models under
every initialisation we tried, so it carries no operational content for LoRA
merging as practised (\S\ref{subsec:floor-zero}); its rate term rests on an
assumption we state but do not prove and on a constant imported from
TurboQuant~\citep{zandieh2025turboquant}, and our attempt to measure the
exponent it predicts has too little dynamic range to do so
(Appendix~\ref{app:rate}); and it constrains a merged model carrying no rank
constraint, while every merge we run is truncated to rank $r$
(\S\ref{subsec:conditioning}). What it did do is put the geometry of the union
of the task subspaces at the centre of the problem, which is what sent us to
measure principal cosines. Nothing else we tried pointed there, and everything
in \S\ref{sec:regimes} and \S\ref{sec:methods} stands or falls on its own
measurements.
